\chapter{LQR/LQG 与 Riccati 方程：存在唯一性、分离原理与 \texorpdfstring{$H_\infty$}{H-infinity} 引入}
\label{ch:lqr-lqg-riccati}

% ════════════════════════════════════════════════════════════════
%  P4 控制部 · C1a 第 C 章「LQR/LQG · Riccati 深化」
%  源：Robotics_Tutorial 01_数学/40_控制理论/50_LQR_LQG与Riccati方程.md（2055 行）
%  定位：深化章。ch:control（导论）已覆盖 LQR 入门全部（DP/HJB 双推、CARE/DARE、
%        Hamilton 矩阵闭式、裕度结论、Doyle 警告、LQG/分离/对偶陈述、P 双关）；
%        C0 五章已做透状态空间/能控能观/Lyapunov/综合/最优控制基础。
%        本章增量 = 把"陈述但未证"的四件事补成教材级证明，再向上接 H∞/数值/学习：
%        (1) CARE 存在唯一性完整证明  (2) 分离原理完整证明
%        (3) 回差不等式证明 + disk margin + LQG 失裕机理  (4) H∞ 博弈 Riccati 引入
%        外加 ARE 数值算法谱系、值/策略迭代=RL 的精确对应。
%  记号归一（设计稿 §3）：Hamilton 矩阵 = Z（区别 Hamilton 函数 H、H∞）；
%        辛矩阵 J=[[0,I],[-I,0]]；共态 lambda；权重 Q,R / 值函数 P；
%        噪声协方差 Σw,Σv / 误差协方差 Σ / Kalman 增益 L。
%  勘误（设计稿 §2，复审清单 C 章 = 1 真错 + 3 部分，全按设计稿 §2 应用）：
%        E1 Doyle 反例用二阶系统（非标量），Q 取 rank-1 [[1,1],[1,1]]；
%        E2 ICL 案例 Garg 勿当 LQR 论文，主引 Goel-Bartlett；
%        E3 "最多 2^n 解"加限定（无虚轴+互异特征值）；
%        E4 Butterworth 限定为 cheap-control / 对称根轨迹语境。
% ════════════════════════════════════════════════════════════════

\begin{practice}[前置自测]
开始之前，先试答下列问题。若已能流畅作答，可只速读本章的几何图解与速查卡；若卡住，正是本章要为你解决的。
\begin{enumerate}
  \item \cref{ch:control} 与 \cref{ch:optimal-control-basics} 已经给出了无限时域 LQR 的最优反馈 $\mathbf u^\star=-\mathbf R^{-1}\mathbf B^\top\mathbf P\mathbf x$ 与连续代数 Riccati 方程（CARE）。但 CARE 是一个\emph{二次}矩阵方程，一般有多个对称解——你能说清，在什么条件下它恰有\emph{一个}使闭环稳定的半正定解？
  \item 都说“LQR 自带 $\ge60^\circ$ 相位裕度”，可一旦把状态换成 Kalman 估计（即 LQG），裕度可以坍缩到零（Doyle 1978）。这两件看似矛盾的事，根子在哪一步？
  \item 分离原理说“状态估计与状态反馈可以分开设计、组合仍最优”。这听上去太好了——它\emph{凭什么}成立？又会在什么时候\emph{失效}？
  \item Riccati 的 Ansatz $\boldsymbol\lambda=\mathbf P\mathbf x$（共态正比于状态）是怎么“猜”出来的？它和 Hamilton 矩阵 $2n$ 个特征值的“稳定一半”有什么几何关系？
  \item 把 LQR 的“平均最优”改成“最坏情况最优”，代价函数与 Riccati 方程会变成什么样？为什么某个临界扰动水平 $\gamma^\star$ 之下，Riccati 就“无解”了？
\end{enumerate}
\end{practice}

\section*{本章目标}
学完本章，你应当能够：
\begin{enumerate}
  \item \textbf{证明}（而不只是陈述）无限时域 LQR 的存在唯一性定理：可稳定化 $+$ 可检测 $\Rightarrow$ CARE 有唯一稳定化半正定解、闭环 Hurwitz；
  \item \textbf{用辛几何}解释 Riccati 的 Ansatz $\boldsymbol\lambda=\mathbf P\mathbf x$ 与 Hamilton 矩阵稳定不变子空间表示 $\mathbf P=\mathbf X_2\mathbf X_1^{-1}$ 的必然性；
  \item \textbf{证明}回差不等式 $\mathbf F^*\mathbf R\mathbf F\succeq\mathbf R$，由此给出 LQR 单输入裕度的频域证明，并理解多输入 disk margin 的退化；
  \item \textbf{完整证明}分离原理（确定性等价），并讲清它在非线性/非高斯/约束下为何失效；
  \item \textbf{剖析} Doyle 1978 反例的机理——Kalman 滤波器作为开环观测器，过程噪声沿不可观方向注入导致裕度坍缩；
  \item \textbf{把} LQR 的 $H_2$ 最优升级为零和微分博弈，推出博弈型 Riccati，并理解 $\gamma\to\infty$ 退化为 CARE、$\gamma\downarrow\gamma^\star$ 时解消失；
  \item \textbf{比较} ARE 的数值算法谱系（Schur / Newton--Kleinman / SDA / 符号函数）及其收敛性；
  \item \textbf{建立} 值迭代/策略迭代 $=$ Riccati 递推/Newton--Kleinman $=$ RL 的精确代数对应，作为通往强化学习的桥。
\end{enumerate}

\begin{note}
\textbf{本章记号与定位.}\;
本章是\emph{深化章}。LQR 的入门级内容——四条路径（HJB / DP / PMP / 完成平方）推 Riccati、CARE/DARE、最优增益、Hamilton 矩阵闭式、裕度结论、LQG/分离/对偶的\emph{陈述}——已在 \cref{ch:control}（\cref{sec:ctrl-lqr-rde,eq:ctrl-care,eq:ctrl-dare}）与 C0 子部 \cref{ch:optimal-control-basics} 中给出，本章\textbf{不重复推导}，只在用到处“回顾一行 $+$ 交叉引用”，把笔墨集中在四个完整证明与向上接口。
记号沿用全书控制部纪律（\cref{tab:ctrl-symbols}）：$\mathbf P$ 是值函数（cost-to-go）权重、$\mathbf Q\succeq\mathbf0,\mathbf R\succ\mathbf0$ 是\emph{代价}权重（均非协方差）；估计侧的过程/观测噪声协方差记 $\boldsymbol\Sigma_w,\boldsymbol\Sigma_v$、误差协方差 $\boldsymbol\Sigma$、Kalman 增益 $\mathbf L$（与 \cref{ch:eskf} 一致）。共态/伴随记 $\boldsymbol\lambda$。
为避免与 Hamilton \emph{函数} $\mathcal H$ 及 $H_\infty$ 的 $H$ 混淆，本章把 Hamilton \textbf{矩阵} 统一记 $\mathcal Z=\begin{psmallmatrix}\mathbf A&-\mathbf B\mathbf R^{-1}\mathbf B^\top\\-\mathbf Q&-\mathbf A^\top\end{psmallmatrix}$，辛矩阵记 $\mathbf J=\begin{psmallmatrix}\mathbf0&\mathbf I\\-\mathbf I&\mathbf0\end{psmallmatrix}$。
\end{note}

% ════════════════════════════════════════════════════════════════
\section{引入与定位：从“能算”到“算得对、算得稳、扛得住”}
\label{sec:lqr-deep-intro}

\paragraph{承上：导论给了“怎么算”，本章补“为什么对”。}
在 \cref{ch:control} 里，我们已经能从动态规划与 HJB 两条路推出 Riccati 方程，写出无限时域 LQR 的最优反馈 $\mathbf u^\star=-\mathbf K^\star\mathbf x$、$\mathbf K^\star=\mathbf R^{-1}\mathbf B^\top\mathbf P^\star$，并知道 $\mathbf P^\star$ 满足 CARE（\cref{eq:ctrl-care}）。\cref{ch:optimal-control-basics} 又从变分与极小值原理把同一结论“第一次系统讲清”。然而这些章节都在“会算”的层面停住了：它们\emph{陈述}了若干关键事实，却把证明留给了本章。具体说，有四件事被反复用到、却尚未证明：

\begin{enumerate}
  \item \textbf{CARE 凭什么有“唯一稳定化解”？} 二次方程一般多解；\cref{thm:ctrl-lqr-converge} 只给了结论。
  \item \textbf{LQR 的好裕度从何而来？} \cref{sec:ctrl-margins} 给了 $\mathrm{GM}\in[\tfrac12,\infty)$、$\mathrm{PM}\ge60^\circ$，但没有证明。
  \item \textbf{分离原理凭什么成立？} \cref{sec:ctrl-lqe-dual} 陈述了“估计与控制可分开设计”，但没有给出正交分解的推导。
  \item \textbf{LQG 为何会丧失裕度？} 导论给了 Doyle 1978 的警告，却没讲清坍缩的机理。
\end{enumerate}

\begin{insight}
本章的增量价值，可一句话概括为：把 \cref{ch:control} 里“\emph{陈述但未证}”的四件事补成教材级证明，再向上接 $H_\infty$、ARE 数值谱系与学习型控制。这与全书控制部的“密度梯度”铁律一致——导论给全景、C0 做透基础、A--Q 研究生密度严格化，同一主题三层各占一档、互不重复。
\end{insight}

\paragraph{LQR 的三层定位。}
为什么值得为一个“线性 $+$ 二次”的特例花一整章？因为 LQR 是最优控制理论的\emph{果蝇}\cite{tedrake_underactuated}：动力学 $\dot{\mathbf x}=\mathbf A\mathbf x+\mathbf B\mathbf u$ 线性、代价二次，结构简单到可以把每条假设都校验一遍，又复杂到涵盖现实工程里大量稳定化任务。它有三重身份：

\begin{itemize}
  \item \textbf{$H_2$ 最优控制}：LQR 的最优代价恰是某闭环传递函数的 $H_2$ 范数平方（\cref{sec:lqr-existence}），由此一步之遥便是 $H_\infty$ 的“最坏情况最优”（\cref{sec:lqr-hinf}）；
  \item \textbf{估计的镜像}：LQR 的控制 Riccati 与 Kalman 滤波的估计 Riccati 互为对偶（\cref{sec:lqr-separation}），一套代码两边通用；
  \item \textbf{强化学习的基准}：值迭代、策略迭代、策略梯度在 LQR 上都有\emph{精确}的代数对应（\cref{sec:lqr-numerics}），Fazel 等称之为“RL 理论的 benchmark”\cite{fazel2018global}。
\end{itemize}

\paragraph{启下：本章脉络与兄弟章依赖。}
本章先把纯 LQ/LQG 的主体讲透（\cref{sec:lqr-finite-track,sec:lqr-existence,sec:lqr-symplectic,sec:lqr-margins,sec:lqr-separation}），再向上打三个接口：$H_\infty$ 的代数入口（\cref{sec:lqr-hinf}，频域范数/$\mu$/DGKF 归兄弟章 \cref{ch:robust-hinf}）、ARE 数值谱系（\cref{sec:lqr-numerics}）、前沿出口（\cref{sec:lqr-frontier}）。依赖关系如下图：变分/PMP（\cref{ch:opt-variational-pmp}）与 DP/HJB（\cref{ch:dp-hjb}）是上游，Lyapunov（\cref{ch:lyapunov-stability}）供稳定性工具，频域 $H_\infty$（\cref{ch:robust-hinf}）与 iLQR/DDP（\cref{ch:ddp-ilqr}）是下游。

\begin{center}
\small
\begin{tikzpicture}[
  node distance=5mm and 8mm, >=Stealth,
  blk/.style={draw, rounded corners, align=center, font=\small, inner sep=3pt, fill=main!6, draw=main!60!black},
  grp/.style={draw, rounded corners, align=center, font=\small, inner sep=3pt, fill=second!7, draw=second!60!black},
  outb/.style={draw, rounded corners, align=center, font=\small, inner sep=3pt, fill=third!9, draw=third!60!black},
  ln/.style={->, thick, gray!55!black}]
\node[blk] (pmp) {变分/PMP\\ \cref{ch:opt-variational-pmp}};
\node[blk, right=of pmp] (dp) {DP/HJB\\ \cref{ch:dp-hjb}};
\node[grp, below=8mm of pmp] (lqr) {LQR/CARE\\ 存在唯一};
\node[grp, right=of lqr] (lqg) {LQG/分离\\ 对偶};
\node[grp, right=of lqg] (hinf) {$H_\infty$ 博弈\\ Riccati（代数侧）};
\node[blk, below=8mm of lqr, fill=structurecolor!8] (lyap) {Lyapunov\\ \cref{ch:lyapunov-stability}};
\node[outb, below=8mm of lqg] (freq) {频域 $H_\infty/\mu$\\ \cref{ch:robust-hinf}};
\node[outb, below=8mm of hinf] (ilqr) {iLQR/DDP\\ \cref{ch:ddp-ilqr}};
\draw[ln] (pmp)--(lqr); \draw[ln] (dp)--(lqr); \draw[ln] (dp)--(lqg);
\draw[ln] (lqr)--(lqg); \draw[ln] (lqg)--(hinf);
\draw[ln] (lyap.north)--(lqr.south);
\draw[ln] (hinf.south)--(freq.north);
\draw[ln] (lqr.south east)--(ilqr.north west);
\end{tikzpicture}
\end{center}

% ════════════════════════════════════════════════════════════════
\section{从有限到无限时域：稳态 Riccati 的诞生与跟踪 LQT}
\label{sec:lqr-finite-track}

\paragraph{回顾（一行带过）。}
有限时域 LQR 的 Riccati 微分方程（RDE）$-\dot{\mathbf P}=\mathbf A^\top\mathbf P+\mathbf P\mathbf A-\mathbf P\mathbf B\mathbf R^{-1}\mathbf B^\top\mathbf P+\mathbf Q$、$\mathbf P(T)=\mathbf Q_f$，及其离散对应（DARE，\cref{eq:ctrl-dare}）、CARE（\cref{eq:ctrl-care}）已在 \cref{sec:ctrl-lqr-rde} 推导，此处不重复。本节只补两件导论未深入的事：有限$\to$无限时域的\emph{收敛速率}，以及把“调节到零”推广为“跟踪参考”的 LQT。

\subsection{有限到无限：指数收敛与 MPC 视界}
\label{sec:lqr-rde-converge}

令 $T\to\infty$、$\mathbf Q_f=\mathbf0$。物理直觉是：若任务没有截止时间，“还剩多少时间”不再影响策略，控制器变为\emph{时不变}，RDE 的解收敛到稳态 $\mathbf P^\star$（CARE 解）。关键的定量问题是\emph{收敛多快}——这直接决定 MPC 需要多长的预测视界。

\begin{theorem}[RDE 指数收敛\cite{anderson2007optimal}]\label{thm:lqr-rde-converge}
在 CARE 良定条件（\cref{thm:lqr-care-exist}）下，RDE 从任意终端 $\mathbf P(T)=\mathbf Q_f\succeq\mathbf0$ 出发，以指数速率收敛到稳态：
\begin{equation}
  \|\mathbf P(t)-\mathbf P^\star\|\le C\,e^{-2\alpha(T-t)},\qquad t\le T,
  \label{eq:lqr-rde-rate}
\end{equation}
其中 $\alpha=\min_i|\operatorname{Re}\lambda_i(\mathbf A_{cl}^\star)|$ 是闭环 $\mathbf A_{cl}^\star=\mathbf A-\mathbf B\mathbf K^\star$ 最慢衰减模态的衰减率。
\end{theorem}

\begin{insight}
\cref{eq:lqr-rde-rate} 本质上是“值迭代”在连续时间的表现：每倒积一步 $\Delta t$ 的 RDE，相当于做一次 Bellman 更新，值函数从终端条件逐步“扩散”到不动点 $\mathbf P^\star$。
工程含义直白：闭环越快（$\alpha$ 越大），$\mathbf P(t)$ 越快收敛——\emph{快}系统只需短视界 MPC，\emph{慢}系统需长视界。取经验门槛 $T\gtrsim 4/\alpha$，则 $\|\mathbf P(0)-\mathbf P^\star\|\lesssim e^{-8}\approx0.03\%\cdot\|\mathbf Q_f-\mathbf P^\star\|$，有限时域与无限时域几乎无别。这正是 \cref{sec:ctrl-mpc-lqr} 中“取终端代价 $V_f=\mathbf x^\top\mathbf P^\star\mathbf x$ 让 MPC 在视界末端无缝接入 LQR”之所以稳妥的根据。
\end{insight}

\subsection{跟踪问题与仿射 LQR（LQT）}
\label{sec:lqr-lqt}

标准 LQR 把状态调节到零，但机器人常需跟踪时变参考 $\mathbf r(t)$（四旋翼跟路径、机械臂跟关节轨迹）。线性二次跟踪器（LQT）把代价改为
\begin{equation}
  J=\tfrac12\!\int_0^T\!\big((\mathbf x-\mathbf r)^\top\mathbf Q(\mathbf x-\mathbf r)+\mathbf u^\top\mathbf R\mathbf u\big)\mathrm dt+\tfrac12(\mathbf x(T)-\mathbf r(T))^\top\mathbf Q_f(\mathbf x(T)-\mathbf r(T)).
  \label{eq:lqt-cost}
\end{equation}
令误差 $\mathbf e=\mathbf x-\mathbf r$，动力学变为 $\dot{\mathbf e}=\mathbf A\mathbf e+\mathbf B\mathbf u+\mathbf d(t)$，其中仿射项 $\mathbf d(t)=\mathbf A\mathbf r(t)-\dot{\mathbf r}(t)$（若 $\mathbf r$ 恰是无控自由响应则 $\mathbf d=\mathbf0$）。关键是值函数由纯二次型升级为\emph{仿射}二次型 $V(\mathbf e,t)=\tfrac12\mathbf e^\top\mathbf P\mathbf e+\mathbf e^\top\mathbf g+c$。代入 HJB 并配方，最优控制取仿射形式
\begin{equation}
  \mathbf u^\star(t)=-\mathbf K(t)(\mathbf x-\mathbf r)-\mathbf R^{-1}\mathbf B^\top\mathbf g(t),
  \label{eq:lqt-u}
\end{equation}
其中 $\mathbf P(t)$ 仍满足\emph{标准} RDE（与参考无关！），而 $\mathbf g(t)$ 满足线性倒向 ODE
\begin{equation}
  -\dot{\mathbf g}=(\mathbf A-\mathbf B\mathbf R^{-1}\mathbf B^\top\mathbf P)^\top\mathbf g+\mathbf P\mathbf d,\qquad \mathbf g(T)=\mathbf Q_f(\mathbf x(T)-\mathbf r(T)).
  \label{eq:lqt-g}
\end{equation}

\begin{insight}
LQT 把控制分解为两部分：\textbf{反馈项} $-\mathbf K(t)(\mathbf x-\mathbf r)$ 抑制偏差（增益由 Riccati 决定）；\textbf{前馈项} $-\mathbf R^{-1}\mathbf B^\top\mathbf g(t)$ 预补偿参考变化（由线性 ODE \cref{eq:lqt-g} 决定）。前馈是 LQT 相对 LQR 的本质新增——没有它，跟踪会有持续稳态误差，尤其对快变参考。自动驾驶横向控制里的“前瞻量”本质上就是这个 $\mathbf g(t)$。
\end{insight}

无穷时域、常值参考 $\mathbf r=\mathrm{const}$ 时，$\mathbf P=\mathbf P^\star$，$\mathbf g$ 满足稳态线性方程 $(\mathbf A-\mathbf B\mathbf K^\star)^\top\mathbf g_\infty=-\mathbf P^\star\mathbf A\mathbf r$；因 $\mathbf A_{cl}^\star$ Hurwitz，其转置亦 Hurwitz，方程有唯一解，稳态前馈 $\mathbf u_{ff}=-\mathbf R^{-1}\mathbf B^\top\mathbf g_\infty$。

\begin{remark}[Bryson 规则作为调参起点]
权重如何选？工程上常以 Bryson 规则起手：取对角权重 $Q_{ii}=1/x_{i,\max}^2$、$R_{jj}=1/u_{j,\max}^2$，即按各状态/输入的“可接受最大值”归一化，使各项在量纲上可比。这只是\emph{起点}而非最优——它把“多大算大”这一物理判断显式化，再据闭环响应迭代微调（参数单调性见 \cref{thm:lqr-monotone}）。
\end{remark}

% ════════════════════════════════════════════════════════════════
\section{CARE 存在唯一性的完整证明}
\label{sec:lqr-existence}

本节是本章第一块核心。\cref{thm:ctrl-lqr-converge} 已\emph{陈述}无限时域 LQR 的存在唯一性，且 \cref{ex:ctrl-kalman} 已用一个“有限时域最优却闭环发散”的反例点明无限时域的必要性；这里把存在唯一性补成完整证明，并梳理 CARE 解集的代数结构。

\subsection{定理陈述与三步证明}
\label{sec:lqr-care-exist}

\begin{theorem}[CARE 存在唯一性，Kalman--Wonham--Anderson--Moore\cite{kalman1960,wonham1968riccati,anderson2007optimal}]\label{thm:lqr-care-exist}
设 $(\mathbf A,\mathbf B)$ 可稳定化、$(\mathbf A,\mathbf Q^{1/2})$ 可检测，$\mathbf R\succ\mathbf0$、$\mathbf Q\succeq\mathbf0$。则 CARE
\begin{equation}
  \mathbf A^\top\mathbf P+\mathbf P\mathbf A-\mathbf P\mathbf B\mathbf R^{-1}\mathbf B^\top\mathbf P+\mathbf Q=\mathbf0
  \label{eq:lqr-care}
\end{equation}
存在\emph{唯一}对称半正定解 $\mathbf P^\star\succeq\mathbf0$，且闭环 $\mathbf A_{cl}^\star=\mathbf A-\mathbf B\mathbf R^{-1}\mathbf B^\top\mathbf P^\star$ 的全部特征值具严格负实部（Hurwitz）。
\end{theorem}

\begin{proof}[证明（三步）]
\textbf{第一步：存在性 $=$ RDE 单调收敛。}\;
取 $\mathbf Q_f=\mathbf0$，从 $t=T$ 倒积 RDE 得 $\mathbf P_T(t)$（下标 $T$ 记终端时刻）。由值函数解释 $\tfrac12\mathbf x^\top\mathbf P_T(0)\mathbf x=\min_{\mathbf u}J_{[0,T]}$，延长时域只增代价（多积一段非负的阶段代价），故 $T\le T'\Rightarrow\mathbf P_T(0)\preceq\mathbf P_{T'}(0)$，即 $\{\mathbf P_T(0)\}_{T>0}$ 单调不减。

\emph{上界}由可稳定化提供：存在 $\mathbf K_0$ 使 $\mathbf A_0=\mathbf A-\mathbf B\mathbf K_0$ Hurwitz。取\emph{次优}控制 $\mathbf u=-\mathbf K_0\mathbf x$，其代价给出上界
\begin{equation}
  \tfrac12\mathbf x^\top\mathbf P_T(0)\mathbf x\le\tfrac12\mathbf x^\top\bar{\mathbf P}\mathbf x,\qquad
  \bar{\mathbf P}:=\int_0^\infty e^{\mathbf A_0^\top t}\big(\mathbf Q+\mathbf K_0^\top\mathbf R\mathbf K_0\big)e^{\mathbf A_0 t}\,\mathrm dt,
  \label{eq:lqr-upper-bound}
\end{equation}
其中 $\bar{\mathbf P}$ 有限（$\mathbf A_0$ Hurwitz $\Rightarrow$ 积分收敛）。于是 $\mathbf0\preceq\mathbf P_T(0)\preceq\bar{\mathbf P}$ 对所有 $T$。单调有界（在半正定锥上）序列收敛，记 $\mathbf P^\star:=\lim_{T\to\infty}\mathbf P_T(0)$；它是 RDE 的稳态（$\dot{\mathbf P}=\mathbf0$），故满足 CARE \cref{eq:lqr-care}。

\textbf{第二步：闭环稳定 $=$ 用可检测性（Lyapunov $+$ LaSalle）。}\;
把 CARE 改写成 Lyapunov 形式：以 $\mathbf K^\star=\mathbf R^{-1}\mathbf B^\top\mathbf P^\star$、$\mathbf A_{cl}^\star=\mathbf A-\mathbf B\mathbf K^\star$，配方得
\begin{equation}
  \mathbf A_{cl}^{\star\top}\mathbf P^\star+\mathbf P^\star\mathbf A_{cl}^\star=-(\mathbf Q+\mathbf K^{\star\top}\mathbf R\mathbf K^\star)\preceq-\mathbf Q.
  \label{eq:lqr-lyap-form}
\end{equation}
取 $V(\mathbf x)=\mathbf x^\top\mathbf P^\star\mathbf x$，沿闭环 $\dot V=-\mathbf x^\top(\mathbf Q+\mathbf K^{\star\top}\mathbf R\mathbf K^\star)\mathbf x\le-\mathbf x^\top\mathbf Q\mathbf x\le0$。若 $\mathbf Q\succ\mathbf0$ 则 $\dot V<0$（$\mathbf x\ne\mathbf0$），直接得 Hurwitz。若仅 $\mathbf Q\succeq\mathbf0$，用 LaSalle：$\dot V=0$ 的集合落在 $\{\mathbf x:\mathbf Q^{1/2}\mathbf x=\mathbf0\}$。在该集合上轨迹满足 $\mathbf Q^{1/2}\mathbf x(t)\equiv\mathbf0$，即状态沿 $\mathbf A_{cl}^\star$ 的\emph{不可观}子空间演化；可检测性 $(\mathbf A,\mathbf Q^{1/2})$ 要求该子空间的所有模态稳定，故 LaSalle 不变集只含原点，$\mathbf A_{cl}^\star$ Hurwitz。（稳定性工具见 \cref{ch:lyapunov-stability} 与 \cref{ch:lyapunov-basics}。）

\textbf{第三步：唯一性 $=$ Lyapunov 方程的唯一性。}\;
设 $\mathbf P'$ 是另一稳定化解（$\mathbf A_{cl}'=\mathbf A-\mathbf B\mathbf R^{-1}\mathbf B^\top\mathbf P'$ 亦 Hurwitz）。令 $\boldsymbol\Delta=\mathbf P^\star-\mathbf P'$，两个 CARE 相减并整理（用 $\mathbf S=\mathbf B\mathbf R^{-1}\mathbf B^\top$）：
\begin{equation}
  \mathbf A_{cl}^{\star\top}\boldsymbol\Delta+\boldsymbol\Delta\mathbf A_{cl}'=\mathbf0
  \quad\Longleftrightarrow\quad
  \mathbf A_{cl}^{\star\top}\boldsymbol\Delta+\boldsymbol\Delta\mathbf A_{cl}^\star=-\boldsymbol\Delta\mathbf S\boldsymbol\Delta\preceq\mathbf0.
  \label{eq:lqr-uniq}
\end{equation}
右式是以 Hurwitz 的 $\mathbf A_{cl}^\star$ 为系数、右端半负定的 Lyapunov 方程，故 $\boldsymbol\Delta\succeq\mathbf0$。对调 $\mathbf P^\star,\mathbf P'$ 的角色（用 $\mathbf A_{cl}'$ Hurwitz）得 $\boldsymbol\Delta\preceq\mathbf0$。两者合并 $\boldsymbol\Delta=\mathbf0$，即 $\mathbf P^\star=\mathbf P'$。
\end{proof}

\begin{insight}
两个条件的工程含义可一句话记住：\textbf{可稳定化}（$(\mathbf A,\mathbf B)$）说“我有能力管住所有发散方向”——否则存在不可控的发散模态、$J=\infty$；\textbf{可检测}（$(\mathbf A,\mathbf Q^{1/2})$）说“我的代价函数关心了所有需要关心的方向”——否则代价“看不见”的模态可能发散、解不唯一或闭环不稳。两者缺一，LQR 要么无解、要么解不唯一。
\end{insight}

\subsection{ARE 的代数结构：单调性、极值解与解的个数}
\label{sec:lqr-are-structure}

存在唯一性之外，CARE 还有一组在调参与数值求解中极有用的代数性质。

\begin{theorem}[参数单调性\cite{anderson2007optimal,lancaster1995algebraic}]\label{thm:lqr-monotone}
在 \cref{thm:lqr-care-exist} 条件下，稳定化解关于权重单调：$\mathbf Q_1\preceq\mathbf Q_2\Rightarrow\mathbf P_1^\star\preceq\mathbf P_2^\star$；$\mathbf R_1\preceq\mathbf R_2\Rightarrow\mathbf P_1^\star\succeq\mathbf P_2^\star$。
\end{theorem}

物理上：增大状态惩罚 $\mathbf Q$ 使系统更“急”回原点，最优代价更大、增益更激进；增大控制权重 $\mathbf R$ 使控制更“贵”，不敢大力控制，状态偏差累积更多、总代价反而更大。这条单调性保证调参不会出现“增大惩罚反而增益减小”的反直觉现象——增大 $Q_{ii}$ 总会使对应状态的增益变大。

\begin{theorem}[稳定化解 $=$ 最大半正定解]\label{thm:lqr-maximal}
CARE 的稳定化解 $\mathbf P^\star$ 是所有对称半正定解中的\emph{最大}者：若 $\mathbf P'\succeq\mathbf0$ 满足 CARE，则 $\mathbf P'\preceq\mathbf P^\star$。
\end{theorem}

直觉是：稳定化解对应的代价从\emph{稳定}轨迹积分而来（有限值），它“看见了最多的未来代价”，故在所有解中最大。反过来，若误取一个非稳定化解 $\mathbf P'\prec\mathbf P^\star$ 作反馈基础，得到的 $\mathbf K'=\mathbf R^{-1}\mathbf B^\top\mathbf P'$ 不稳定化系统——闭环在某些方向发散，而 $\mathbf P'$ 没“预见”到这些方向的无穷代价，所以它比 $\mathbf P^\star$ 小。这正是数值求解必须锁定\emph{稳定}不变子空间的原因（\cref{sec:lqr-symplectic,sec:lqr-numerics}）。

\begin{pitfall}[口径陷阱：“CARE 一般有 $2^n$ 个解”要加限定]
常见的不严谨说法是“CARE 最多有 $2^n$ 个对称解”。精确陈述是：实对称解与 Hamilton 矩阵 $\mathcal Z$ 的\emph{Lagrangian 且共轭闭}不变子空间一一对应。\textbf{当 $\mathcal Z$ 的 $2n$ 个特征值两两互异、且无虚轴特征值时}，按“每对共轭特征值分到稳定侧或反稳定侧”最多有 $2^n$ 种取法、对应至多 $2^n$ 个解；有重特征值时解构成连续族，有虚轴特征值时良定性本身被破坏\cite{lancaster1995algebraic}。在标准 LQR 中我们只取“全部稳定特征值”那一支，即 $\mathbf P^\star$；其余子空间在 $H_\infty$ 博弈（\cref{sec:lqr-hinf}）中才派上用场。
\end{pitfall}

\begin{table}[H]\centering\small
\caption{三类矩阵方程的层级：从耦合到最优}
\label{tab:lqr-ric-lyap}
\begin{tabular}{lll}
\toprule
方程 & 形式 & 含义 \\
\midrule
Sylvester & $\mathbf A\mathbf X+\mathbf X\mathbf B=\mathbf C$（线性） & 两系统间的耦合 \\
Lyapunov & $\mathbf A^\top\mathbf P+\mathbf P\mathbf A=-\mathbf Q$（线性，Sylvester 特例） & 给定反馈的闭环代价 \\
Riccati & $\mathbf A^\top\mathbf P+\mathbf P\mathbf A-\mathbf P\mathbf B\mathbf R^{-1}\mathbf B^\top\mathbf P+\mathbf Q=\mathbf0$（二次） & \emph{最优}反馈的闭环代价 \\
\bottomrule
\end{tabular}
\end{table}

\cref{tab:lqr-ric-lyap} 揭示了一条关键退化：固定某反馈 $\mathbf K$（不做优化），CARE 退化为 Lyapunov 方程 $\mathbf A_K^\top\mathbf P_K+\mathbf P_K\mathbf A_K+\mathbf Q+\mathbf K^\top\mathbf R\mathbf K=\mathbf0$（$\mathbf A_K=\mathbf A-\mathbf B\mathbf K$）。Riccati 比 Lyapunov 多出的 $-\mathbf P\mathbf B\mathbf R^{-1}\mathbf B^\top\mathbf P$ 项，正是“优化”的代价。Newton--Kleinman 迭代（\cref{sec:lqr-numerics}）就利用这一层级：在当前 $\mathbf K_k$ 处解线性 Lyapunov（快），再更新 $\mathbf K_{k+1}$。

% ════════════════════════════════════════════════════════════════
\section{PMP 视角与辛几何：Hamilton 矩阵、Lagrangian 子空间}
\label{sec:lqr-symplectic}

\cref{eq:ctrl-hamilton-sol} 给了用 Hamilton 矩阵解 CARE 的\emph{闭式}（“how”）。本节回答“why”：Riccati 的 Ansatz $\boldsymbol\lambda=\mathbf P\mathbf x$ 不是天才的猜测，而是辛几何的必然。PMP 的基本五步推导见 \cref{ch:opt-variational-pmp}，此处只回顾接口。

\paragraph{从 PMP 到 Hamilton 矩阵。}
极小值原理（\cref{ch:opt-variational-pmp}）对 LQR 给出共态方程 $\dot{\boldsymbol\lambda}=-\partial\mathcal H/\partial\mathbf x=-\mathbf A^\top\boldsymbol\lambda-\mathbf Q\mathbf x$ 与 $\mathbf u^\star=-\mathbf R^{-1}\mathbf B^\top\boldsymbol\lambda$。把状态、共态合成 $2n$ 维线性系统：
\begin{equation}
  \begin{pmatrix}\dot{\mathbf x}\\\dot{\boldsymbol\lambda}\end{pmatrix}
  =\mathcal Z\begin{pmatrix}\mathbf x\\\boldsymbol\lambda\end{pmatrix},\qquad
  \mathcal Z=\begin{pmatrix}\mathbf A&-\mathbf S\\-\mathbf Q&-\mathbf A^\top\end{pmatrix},\quad \mathbf S:=\mathbf B\mathbf R^{-1}\mathbf B^\top\succeq\mathbf0.
  \label{eq:lqr-Z}
\end{equation}

\begin{definition}[Hamilton 矩阵与辛矩阵]\label{def:lqr-hamiltonian}
设辛矩阵 $\mathbf J=\begin{psmallmatrix}\mathbf0&\mathbf I\\-\mathbf I&\mathbf0\end{psmallmatrix}\in\mathbb R^{2n\times2n}$。称 $\mathcal Z\in\mathbb R^{2n\times2n}$ 为 \emph{Hamilton 矩阵}（属辛李代数 $\mathfrak{sp}(2n)$），若 $\mathbf J\mathcal Z+\mathcal Z^\top\mathbf J=\mathbf0$，等价于 $\mathbf J\mathcal Z$ 对称。
\end{definition}

\begin{proposition}[LQR 的 $\mathcal Z$ 是 Hamilton 矩阵，且谱关于虚轴对称]\label{prop:lqr-symplectic}
\cref{eq:lqr-Z} 的 $\mathcal Z$ 满足 \cref{def:lqr-hamiltonian}；其特征值关于虚轴对称——若 $\mu\in\sigma(\mathcal Z)$ 则 $-\bar\mu\in\sigma(\mathcal Z)$（实矩阵时简化为 $-\mu\in\sigma(\mathcal Z)$）。在 CARE 良定条件下 $\mathcal Z$ 无虚轴特征值，故恰有 $n$ 个特征值在开左半平面、$n$ 个在开右半平面。
\end{proposition}

\begin{proof}
$\mathbf J\mathcal Z=\begin{psmallmatrix}\mathbf0&\mathbf I\\-\mathbf I&\mathbf0\end{psmallmatrix}\begin{psmallmatrix}\mathbf A&-\mathbf S\\-\mathbf Q&-\mathbf A^\top\end{psmallmatrix}=\begin{psmallmatrix}-\mathbf Q&-\mathbf A^\top\\-\mathbf A&\mathbf S\end{psmallmatrix}$，对称（$\mathbf Q,\mathbf S$ 对称、非对角块互为转置），故 $\mathcal Z$ 是 Hamilton 矩阵。由 $(\mathbf J\mathcal Z)^\top=\mathbf J\mathcal Z$ 得 $\mathcal Z^\top\mathbf J^\top=\mathbf J\mathcal Z$，即 $-\mathcal Z^\top=\mathbf J\mathcal Z\mathbf J^{-1}$，故 $\mathcal Z$ 与 $-\mathcal Z^\top$ 相似、同谱；而 $-\mathcal Z^\top$ 的谱是 $\sigma(\mathcal Z)$ 取负再共轭，即 $\{-\bar\mu\}$。无虚轴特征值的结论由 $\mu$ 与 $-\bar\mu$ 不可同时在虚轴外又落到虚轴上、配合可稳定化 $+$ 可检测排除虚轴解得到。
\end{proof}

\paragraph{稳定不变子空间表示。}
设 $\mathcal Z$ 的 $n$ 个稳定特征值对应的稳定不变子空间为 $\mathcal V_-\subset\mathbb R^{2n}$（$\dim=n$），把其一组基写成 $\begin{psmallmatrix}\mathbf X_1\\\mathbf X_2\end{psmallmatrix}$，$\mathbf X_1,\mathbf X_2\in\mathbb R^{n\times n}$。

\begin{theorem}[稳定不变子空间表示\cite{anderson2007optimal,laub1979schur}]\label{thm:lqr-invsubspace}
若 $\mathbf X_1$ 可逆，则 $\mathbf P^\star=\mathbf X_2\mathbf X_1^{-1}$ 是 CARE 的唯一稳定化解。
\end{theorem}

\begin{proof}
不变子空间满足 $\mathcal Z\begin{psmallmatrix}\mathbf X_1\\\mathbf X_2\end{psmallmatrix}=\begin{psmallmatrix}\mathbf X_1\\\mathbf X_2\end{psmallmatrix}\boldsymbol\Lambda$，$\sigma(\boldsymbol\Lambda)\subset\mathbb C_-$。展开 \cref{eq:lqr-Z}：$\mathbf A\mathbf X_1-\mathbf S\mathbf X_2=\mathbf X_1\boldsymbol\Lambda$、$-\mathbf Q\mathbf X_1-\mathbf A^\top\mathbf X_2=\mathbf X_2\boldsymbol\Lambda$。令 $\mathbf P=\mathbf X_2\mathbf X_1^{-1}$。由第一式右乘 $\mathbf X_1^{-1}$ 得 $\mathbf A-\mathbf S\mathbf P=\mathbf X_1\boldsymbol\Lambda\mathbf X_1^{-1}$，故闭环 $\mathbf A-\mathbf B\mathbf K$ 与 $\boldsymbol\Lambda$ 相似、Hurwitz。第二式右乘 $\mathbf X_1^{-1}$ 得 $-\mathbf Q-\mathbf A^\top\mathbf P=\mathbf P\mathbf X_1\boldsymbol\Lambda\mathbf X_1^{-1}=\mathbf P(\mathbf A-\mathbf S\mathbf P)$，整理即 $\mathbf A^\top\mathbf P+\mathbf P\mathbf A-\mathbf P\mathbf S\mathbf P+\mathbf Q=\mathbf0$，正是 CARE。
\end{proof}

\begin{insight}
\cref{thm:lqr-invsubspace} 揭穿了 Riccati Ansatz 的“魔术”。$2n$ 维相空间 $(\mathbf x,\boldsymbol\lambda)$ 配上辛形式 $\mathbf J$，稳定不变子空间 $\mathcal V_-$ 恰是一个 $n$ 维 \emph{Lagrangian} 子空间（即在 $\mathbf J$ 下自正交）。而任何“竖直分量 $\mathbf X_1$ 可逆”的 Lagrangian 子空间都能唯一地写成一张\emph{图}：$\mathcal V_-=\{(\mathbf x,\mathbf P\mathbf x):\mathbf x\in\mathbb R^n\}$，其中 $\mathbf P=\mathbf P^\top$。换言之，$\boldsymbol\lambda=\mathbf P\mathbf x$ 不是猜出来的——它是“稳定 Lagrangian 子空间 $=$ 某对称 $\mathbf P$ 的图”这一辛几何事实的改写，而 Riccati 方程正是该子空间在 Hamilton 流下保持不变的条件。$\mathbf P$ 的对称性，对应 Lagrangian 性质。
\end{insight}

\begin{remark}[与导论的分工]
\cref{eq:ctrl-hamilton-sol} 给“$2n\times2n$ Hamilton 矩阵 $+$ 稳定不变子空间”的闭式算法（“how”），本节给其辛几何根据（“why”）。数值实现（实 Schur $+$ 特征值排序提取 $\mathbf P=\mathbf U_{21}\mathbf U_{11}^{-1}$）留到 \cref{sec:lqr-numerics}。
\end{remark}

% ════════════════════════════════════════════════════════════════
\section{LQR 鲁棒裕度的证明，及 LQG 为何丧失它}
\label{sec:lqr-margins}

本节是第二块核心。\cref{sec:ctrl-margins} 已陈述 LQR 单输入裕度 $\mathrm{GM}\in[\tfrac12,\infty)$、$\mathrm{PM}\ge60^\circ$，并警告 LQG 会丧失裕度。这里给出回差不等式的证明、多输入 disk margin 的退化，以及 Doyle 反例的机理。

\subsection{回差不等式与单输入裕度的证明}
\label{sec:lqr-return-diff}

\begin{definition}[回差函数]\label{def:lqr-return-diff}
LQR 的\emph{回差}（return difference）为 $\mathbf F(s)=\mathbf I+\mathbf K(s\mathbf I-\mathbf A)^{-1}\mathbf B$，$\mathbf K=\mathbf R^{-1}\mathbf B^\top\mathbf P^\star$。$\mathbf F(j\omega)$ 衡量“闭环相对开环的改善”，开环传递为 $\mathbf L(s)=\mathbf K(s\mathbf I-\mathbf A)^{-1}\mathbf B$。
\end{definition}

\begin{theorem}[回差不等式，Kalman 1964\cite{kalman1964optimal}]\label{thm:lqr-kalman-ineq}
对 LQR 最优反馈，回差满足
\begin{equation}
  \mathbf F(j\omega)^*\mathbf R\,\mathbf F(j\omega)=\mathbf R+\mathbf B^\top(-j\omega\mathbf I-\mathbf A^\top)^{-1}\mathbf Q(j\omega\mathbf I-\mathbf A)^{-1}\mathbf B\succeq\mathbf R,\qquad\forall\omega\in\mathbb R.
  \label{eq:lqr-return-ineq}
\end{equation}
\end{theorem}

\begin{proof}[证明骨架]
从 Lyapunov 形式的 CARE \cref{eq:lqr-lyap-form}：$\mathbf A_{cl}^{\star\top}\mathbf P^\star+\mathbf P^\star\mathbf A_{cl}^\star+\mathbf Q+\mathbf K^\top\mathbf R\mathbf K=\mathbf0$。在两侧分别左乘 $\mathbf B^\top(-j\omega\mathbf I-\mathbf A^\top)^{-1}$、右乘 $(j\omega\mathbf I-\mathbf A)^{-1}\mathbf B$，并利用恒等式 $(j\omega\mathbf I-\mathbf A)^{-1}-(j\omega\mathbf I-\mathbf A_{cl}^\star)^{-1}=(j\omega\mathbf I-\mathbf A_{cl}^\star)^{-1}\mathbf B\mathbf K(j\omega\mathbf I-\mathbf A)^{-1}$ 把 $\mathbf A_{cl}^\star$ 的分式换成 $\mathbf A$ 的分式与 $\mathbf L(j\omega)=\mathbf K(j\omega\mathbf I-\mathbf A)^{-1}\mathbf B$。代入 $\mathbf K=\mathbf R^{-1}\mathbf B^\top\mathbf P^\star$ 后，左端整理为 $(\mathbf I+\mathbf L)^*\mathbf R(\mathbf I+\mathbf L)=\mathbf F^*\mathbf R\mathbf F$，右端剩下 $\mathbf R$ 加上一项 $\mathbf B^\top(-j\omega\mathbf I-\mathbf A^\top)^{-1}\mathbf Q(j\omega\mathbf I-\mathbf A)^{-1}\mathbf B\succeq\mathbf0$（因 $\mathbf Q\succeq\mathbf0$），即 \cref{eq:lqr-return-ineq}。完整代数见 Anderson--Moore Ch.11\cite{anderson2007optimal} 或 \cref{ch:robust-hinf}。
\end{proof}

\begin{theorem}[单输入裕度，频域证明]\label{thm:lqr-simargin}
单输入（$m=1$、$\mathbf R=r>0$）时 \cref{eq:lqr-return-ineq} 化为 $|1+L(j\omega)|\ge1\ \forall\omega$，即 Nyquist 曲线 $L(j\omega)$ 永不进入以 $-1$ 为圆心、半径 $1$ 的开圆盘。由此 $\mathrm{GM}\in[\tfrac12,\infty)$（$\ge6$\,dB、无上界）、$\mathrm{PM}\ge60^\circ$。
\end{theorem}

\begin{proof}
$|1+L(j\omega)|\ge1$ 即 $L$ 不入 $|s+1|<1$ 的圆盘。\emph{增益裕度}：把 $L$ 乘 $k>0$，只要轨迹不穿 $-1$ 即闭环稳定；圆盘把 $-1$ 右侧到原点的临界段“挡住”，可证 $L$ 可乘任意 $k\in[\tfrac12,\infty)$ 而不穿 $-1$（$k=\tfrac12$ 对应把临界点从 $-2$ 推到 $-1$ 的边界）。\emph{相位裕度}：$|1+L(j\omega)|=1$ 对应 $L$ 落在以 $-1$、$0$ 为端点的单位圆上；当 $|L(j\omega_{gc})|=1$ 时，几何上 $L(j\omega_{gc})$ 与 $-1$ 的张角至少 $60^\circ$，故 $\mathrm{PM}\ge60^\circ$。
\end{proof}

\begin{insight}
这条裕度是\emph{最优性的免费午餐}：回差不等式是 Riccati 解的结构后果，非最优反馈（如手调 PID）没有这一保证。
\textbf{反事实}：若只“近似”解了 ARE（如截断的 Newton--Kleinman，\cref{sec:lqr-numerics}），回差不等式\emph{不一定}成立，次优解的裕度可能差很多——这是工程中必须用高精度求解器（收敛到机器精度）而非粗略近似的理由。
\end{insight}

\subsection{多输入回差与 disk margin}
\label{sec:lqr-disk-margin}

对多输入（$m>1$），Safonov--Athans 1977\cite{safonov1977gain} 证明矩阵版回差不等式 \cref{eq:lqr-return-ineq} 仍成立（$\mathbf F$ 为 $m\times m$）。但它\emph{不}直接给出每个通道的标量裕度：多输入 LQR 对\emph{同时}在所有通道施加相同增益/相位扰动裕度很好，但对\emph{单通道}扰动（仅一个执行器故障）裕度可能很差。现代鲁棒控制用 \emph{disk margin} 取代经典 GM/PM——设通道增益扰动 $\delta_i$ 落在复圆盘 $|\delta_i-1|<\rho$ 内，系统能容忍的最大 $\rho$ 即 disk margin。单输入 LQR 的 disk margin $\ge\tfrac12$（对应经典结果）；多输入则依 $\mathbf B$ 的结构，一般 $<\tfrac12$。

\begin{pitfall}[LQR 裕度的三个前提]
LQR 的裕度保证仅在三前提下有效，逾越即失效：
\begin{enumerate}
  \item \textbf{状态完全可测}：若状态是估计值（LQG），裕度丧失（\cref{sec:lqr-doyle}，Doyle 1978）；
  \item \textbf{模型精确}：裕度是对\emph{回路增益}扰动的，\emph{不是}对\emph{模型参数}扰动的——质量变化 $20\%$ 不一定落在裕度内；
  \item \textbf{无饱和}：执行器一旦饱和，反馈不再线性，裕度失效（饱和处理 $\to$ 约束 MPC，\cref{ch:mpc-advanced}）。
\end{enumerate}
这正是现代实践中 LQR 常作“标称设计”、再用 $H_\infty$/鲁棒 MPC 处理不确定性的原因。
\end{pitfall}

\subsection{LQG 丧失裕度：Doyle 1978 的机理}
\label{sec:lqr-doyle}

分离原理（\cref{sec:lqr-separation}）保证 LQG 性能最优（均方代价），但\emph{不}保证鲁棒裕度。LQR 自带的 $\ge60^\circ$/$[\tfrac12,\infty)$ 裕度，在 LQG 中可坍缩到任意小。

\begin{example}[Doyle 1978 反例\cite{doyle1978guaranteed}]\label{ex:lqr-doyle}
考虑二阶 SISO 系统
\begin{equation}
  \mathbf A=\begin{pmatrix}1&1\\0&1\end{pmatrix},\quad \mathbf B=\begin{pmatrix}0\\1\end{pmatrix},\quad \mathbf C=\begin{pmatrix}1&0\end{pmatrix},
  \label{eq:lqr-doyle-sys}
\end{equation}
开环在 $+1$ 有二重特征值（不稳定）。取代价权重 $\mathbf Q=q\begin{psmallmatrix}1&1\\1&1\end{psmallmatrix}$（\textbf{rank-1、非对角}）、$\mathbf R=1$，过程噪声协方差 $\boldsymbol\Sigma_w=\sigma\begin{psmallmatrix}1&1\\1&1\end{psmallmatrix}$、$\boldsymbol\Sigma_v=1$。当 $q,\sigma\to\infty$ 时，LQG 闭环的增益裕度可\emph{任意趋于 $0$}、相位裕度趋于零。Doyle 论文摘要仅四字——“\emph{There are none.}”——意即 LQG 没有保证的稳定裕度，成为控制史上最短最震撼的摘要之一。
\end{example}

\begin{insight}[坍缩机理]
为什么是\emph{二阶}、为什么 $\mathbf Q$ 取 rank-1 的 $\begin{psmallmatrix}1&1\\1&1\end{psmallmatrix}$（而非 $q\mathbf I$）？因为机理是“\textbf{过程噪声沿不可观/弱可观方向注入}”：rank-1 的 $\boldsymbol\Sigma_w$ 把噪声集中打到一个特定方向，迫使 Kalman 滤波器在该方向极不信任模型、增益 $\mathbf L$ 爆炸；而高增益观测器把极点推深，却在回路断开点引入额外相位滞后，把 LQR 原有的回差性质 $|1+L|\ge1$ 破坏掉。标量一阶系统只有一个方向，无法制造“可观与不可观错位”，故\emph{无法}复现裕度坍缩——这就是反例必须是二阶、$\mathbf Q$ 必须非对角的根本原因。
物理根源一句话：\textbf{Kalman 滤波器是“开环观测器”}——基于标称 $(\mathbf A,\mathbf C)$ 构建，遇参数失配无内在修正机制；LQR 的鲁棒性来自反馈回差，而回差在串入观测器后被观测器极点破坏。
\end{insight}

三条解药顺势引出本章后续与兄弟章：(i) \textbf{LTR}（\cref{sec:lqr-hinf} 末）以虚拟过程噪声恢复 LQR 回路；(ii) \textbf{$H_\infty$ 综合}（\cref{sec:lqr-hinf}）把 $H_2$ 升级为 minimax 博弈、直接对抗最坏扰动；(iii) \textbf{$\mu$-综合}处理\emph{结构化}不确定性——当不确定性具块对角结构时 $H_\infty$ 范数过保守，结构奇异值 $\mu$ 给更紧的度量，DK 迭代交替优化，详见频域章 \cref{ch:robust-hinf}。

% ════════════════════════════════════════════════════════════════
\section{LQG、分离原理完整证明与 Kalman--LQR 对偶}
\label{sec:lqr-separation}

本节是第三块核心。\cref{sec:ctrl-lqe-dual} 已陈述分离原理与 LQR--Kalman 对偶，这里把分离原理补成完整证明，并讲清失效条件。

\subsection{LQG 问题与 Kalman--Bucy 滤波 Riccati}
\label{sec:lqr-lqg-setup}

LQG 问题在 LQR 上叠加了过程噪声与不完全观测：
\begin{equation}
  \dot{\mathbf x}=\mathbf A\mathbf x+\mathbf B\mathbf u+\mathbf w,\quad \mathbf y=\mathbf C\mathbf x+\mathbf v,\quad \mathbf w\sim\mathcal N(\mathbf0,\boldsymbol\Sigma_w),\ \mathbf v\sim\mathcal N(\mathbf0,\boldsymbol\Sigma_v),
  \label{eq:lqr-lqg-plant}
\end{equation}
最小化 $\mathbb E\!\int(\mathbf x^\top\mathbf Q\mathbf x+\mathbf u^\top\mathbf R\mathbf u)\mathrm dt$。未知的是状态 $\mathbf x$，只能观测 $\mathbf y$。Kalman--Bucy 滤波器给出条件均值 $\hat{\mathbf x}=\mathbb E[\mathbf x\mid\mathcal Y_t]$，其动力学与误差协方差 $\boldsymbol\Sigma=\mathbb E[\tilde{\mathbf x}\tilde{\mathbf x}^\top]$（$\tilde{\mathbf x}=\mathbf x-\hat{\mathbf x}$）满足
\begin{equation}
  \dot{\hat{\mathbf x}}=\mathbf A\hat{\mathbf x}+\mathbf B\mathbf u+\mathbf L(\mathbf y-\mathbf C\hat{\mathbf x}),\qquad
  \dot{\boldsymbol\Sigma}=\mathbf A\boldsymbol\Sigma+\boldsymbol\Sigma\mathbf A^\top-\boldsymbol\Sigma\mathbf C^\top\boldsymbol\Sigma_v^{-1}\mathbf C\boldsymbol\Sigma+\boldsymbol\Sigma_w,\quad \mathbf L=\boldsymbol\Sigma\mathbf C^\top\boldsymbol\Sigma_v^{-1}.
  \label{eq:lqr-kbf}
\end{equation}
（连续 Kalman 滤波的完整推导见 \cref{ch:eskf}；此处只强调其作为 LQR 对偶的视角。）

\subsection{分离原理的完整证明}
\label{sec:lqr-sep-proof}

\begin{theorem}[分离原理 / 确定性等价\cite{kwakernaak1972linear,anderson2007optimal}]\label{thm:lqr-separation}
对线性--高斯--二次的 LQG 问题，最优控制器 $=$ Kalman 滤波器 $\hat{\mathbf x}$（\cref{eq:lqr-kbf}）$+$ LQR 反馈 $\mathbf u=-\mathbf K\hat{\mathbf x}$，其中 $\mathbf K=\mathbf R^{-1}\mathbf B^\top\mathbf P^\star$ 是\emph{假设状态可测}时的 LQR 增益。两个子问题可\emph{独立}设计，组合后仍全局最优。
\end{theorem}

\begin{proof}[证明（确定性等价，五步）]
\textbf{Step 1：正交分解。} 估计误差 $\tilde{\mathbf x}=\mathbf x-\hat{\mathbf x}$，$\hat{\mathbf x}=\mathbb E[\mathbf x\mid\mathcal Y_t]$。由正交投影定理，$\tilde{\mathbf x}\perp\hat{\mathbf x}$，即 $\mathbb E[\tilde{\mathbf x}\hat{\mathbf x}^\top]=\mathbf0$。

\textbf{Step 2：代价正交分解。} 状态二次型分解：
\begin{equation}
  \mathbb E[\mathbf x^\top\mathbf Q\mathbf x]=\mathbb E[(\hat{\mathbf x}+\tilde{\mathbf x})^\top\mathbf Q(\hat{\mathbf x}+\tilde{\mathbf x})]=\mathbb E[\hat{\mathbf x}^\top\mathbf Q\hat{\mathbf x}]+\mathbb E[\tilde{\mathbf x}^\top\mathbf Q\tilde{\mathbf x}]+2\underbrace{\mathbb E[\hat{\mathbf x}^\top\mathbf Q\tilde{\mathbf x}]}_{=\,0}.
  \label{eq:lqr-cost-split}
\end{equation}
故总代价 $J=\mathbb E\!\int(\hat{\mathbf x}^\top\mathbf Q\hat{\mathbf x}+\mathbf u^\top\mathbf R\mathbf u)\mathrm dt+\mathbb E\!\int\tilde{\mathbf x}^\top\mathbf Q\tilde{\mathbf x}\,\mathrm dt$。

\textbf{Step 3：估计误差独立于控制。} 对线性--高斯系统，误差协方差 Riccati \cref{eq:lqr-kbf} \emph{不含} $\mathbf u$——控制输入不影响估计精度（这是线性--高斯假设的关键后果）。故 $J$ 的第二项是与 $\mathbf u$ 无关的常数，优化只需处理第一项。

\textbf{Step 4：对 $\hat{\mathbf x}$ 的确定性等价。} $\hat{\mathbf x}$ 的动力学 \cref{eq:lqr-kbf} 中，新息项 $\mathbf L(\mathbf y-\mathbf C\hat{\mathbf x})$ 是零均值白噪声（新息过程），其对期望代价贡献为常数。于是第一项恰是以 $\hat{\mathbf x}$ 为状态、$\mathbf u$ 为控制的\emph{确定性} LQR 问题。

\textbf{Step 5：结论。} 该确定性 LQR 的最优解为 $\mathbf u=-\mathbf K\hat{\mathbf x}$、$\mathbf K=\mathbf R^{-1}\mathbf B^\top\mathbf P^\star$，而 $\hat{\mathbf x}$ 由 Kalman 滤波器产生。两个子问题独立设计、组合后全局最优。
\end{proof}

\begin{insight}
分离原理的精妙在于把\emph{部分可观}的随机最优控制拆成两个\emph{完全可解}的子问题。它\emph{只}在线性--高斯--二次下严格成立——\cref{eq:lqr-cost-split} 的正交项为零、\cref{eq:lqr-kbf} 的误差 Riccati 不含 $\mathbf u$，缺一不可。
\end{insight}

\begin{table}[H]\centering\small
\caption{分离原理的失效条件}
\label{tab:lqr-sep-fail}
\begin{tabular}{p{0.24\linewidth}p{0.70\linewidth}}
\toprule
破坏的假设 & 如何破坏正交分解 \\
\midrule
非线性动力学 & $\mathbb E[\tilde{\mathbf x}\hat{\mathbf x}^\top]\ne\mathbf0$（估计误差与状态估计相关），代价无法正交分解；误差 Riccati 一般含 $\mathbf u$。 \\
非高斯噪声 & 条件分布不由均值/协方差完全表征，最优估计不是 Kalman 滤波器。 \\
非二次代价 & \cref{eq:lqr-cost-split} 的交叉项不再消失。 \\
约束 / 饱和 & 反馈非线性，最优控制不再是 $-\mathbf K\hat{\mathbf x}$。 \\
\bottomrule
\end{tabular}
\end{table}

机器人含义：SLAM 中常用的 EKF $+$ LQR 是\emph{近似}分离——线性化误差小时有效，剧烈运动下 $\mathbb E[\tilde{\mathbf x}\hat{\mathbf x}^\top]\ne\mathbf0$，性能可能意外崩溃（应对手段指向力控/全身控制 WBC 与约束 MPC，见 \cref{ch:planning} 与 \cref{ch:force-wbc}）。

\subsection{Kalman--LQR 对偶}
\label{sec:lqr-duality}

\cref{eq:lqr-kbf} 的滤波 Riccati 与 CARE 结构相同，互为对偶。

\begin{table}[H]\centering\small
\caption{Kalman--LQR 对偶：控制即反向估计}
\label{tab:lqr-duality}
\begin{tabular}{ll}
\toprule
控制（LQR） & 估计（Kalman） \\
\midrule
$\mathbf A$ & $\mathbf A^\top$ \\
$\mathbf B$ & $\mathbf C^\top$ \\
$\mathbf Q$ & $\boldsymbol\Sigma_w$ \\
$\mathbf R$ & $\boldsymbol\Sigma_v$ \\
$\mathbf P^\star$（值函数） & $\boldsymbol\Sigma^\star$（误差协方差） \\
$\mathbf K=\mathbf R^{-1}\mathbf B^\top\mathbf P^\star$ & $\mathbf L=\boldsymbol\Sigma^\star\mathbf C^\top\boldsymbol\Sigma_v^{-1}$ \\
可稳定化 $(\mathbf A,\mathbf B)$ & 可检测 $(\mathbf A,\mathbf C)$ \\
\bottomrule
\end{tabular}
\end{table}

\begin{proposition}[严格对偶]\label{prop:lqr-duality}
Kalman--Bucy 滤波的误差协方差 Riccati \cref{eq:lqr-kbf} 与\emph{对偶系统} $(\mathbf A^\top,\mathbf C^\top,\boldsymbol\Sigma_w,\boldsymbol\Sigma_v)$ 的控制 CARE 是\emph{同一个方程}：在 \cref{eq:lqr-care} 中作替换 $\mathbf A\mapsto\mathbf A^\top$、$\mathbf B\mapsto\mathbf C^\top$、$\mathbf Q\mapsto\boldsymbol\Sigma_w$、$\mathbf R\mapsto\boldsymbol\Sigma_v$，所得 CARE
\[
  \mathbf A\boldsymbol\Sigma+\boldsymbol\Sigma\mathbf A^\top-\boldsymbol\Sigma\mathbf C^\top\boldsymbol\Sigma_v^{-1}\mathbf C\boldsymbol\Sigma+\boldsymbol\Sigma_w=\mathbf0
\]
恰是 \cref{eq:lqr-kbf} 的稳态形式。\footnote{故对偶是“同一方程换一套参数”，而非两组无关参数的解恰好相等：必须按上述字典 $(\mathbf B\!\leftrightarrow\!\mathbf C^\top,\ \mathbf Q\!\leftrightarrow\!\boldsymbol\Sigma_w,\ \mathbf R\!\leftrightarrow\!\boldsymbol\Sigma_v)$ 对应替换，控制 CARE 的可稳定化/可检测前提相应地变成滤波侧的可检测 $(\mathbf A,\mathbf C)$/可稳定化 $(\mathbf A,\boldsymbol\Sigma_w^{1/2})$。}于是 $(\mathbf A,\mathbf B,\mathbf Q,\mathbf R)$ 的控制 Riccati 与 $(\mathbf A,\mathbf C,\boldsymbol\Sigma_w,\boldsymbol\Sigma_v)$ 的滤波 Riccati 可由同一求解器、同一套理论一并处理。
\end{proposition}

\begin{insight}
对偶不是符号巧合，而是“控制即反向估计”的数学表述：控制从现在向未来\emph{注入}信息（经 $\mathbf u$），估计从过去向现在\emph{提取}信息（经 $\mathbf y$）；两者都是沿时间轴的信息传播，方向相反。时间反转 $t\mapsto-t$ 把一个变成另一个。工程后果：同一个 Schur 算法（\cref{sec:lqr-numerics}）既解控制 Riccati 又解滤波 Riccati，换一下矩阵参数即可。这与 \cref{sec:ctrl-lqe-dual} 的“前向 vs 后向、协方差 vs 值函数”呼应——注意 $\mathbf P$ 的双关已在 \cref{tab:ctrl-symbols} 处声明，本章不再重复整段。
\end{insight}

% ════════════════════════════════════════════════════════════════
\section{\texorpdfstring{$H_\infty$}{H-infinity} 引入：从 \texorpdfstring{$H_2$}{H2} 到零和微分博弈 Riccati}
\label{sec:lqr-hinf}

本节是本章上限，也是向鲁棒控制的接口。它只做 $H_\infty$ 的\emph{代数/博弈}入口——状态空间博弈 Riccati 与它和 LQR 的代数连续性；\textbf{频域 $H_\infty$ 范数、灵敏度 $S/T$、$\mu$-综合/DK 迭代、DGKF 解全部归兄弟章 \cref{ch:robust-hinf}}。

\subsection{LQR 即 \texorpdfstring{$H_2$}{H2} 最优}
\label{sec:lqr-h2}

把 LQR 写成传递函数语言：令评价输出 $\mathbf z=\begin{psmallmatrix}\mathbf Q^{1/2}\mathbf x\\\mathbf R^{1/2}\mathbf u\end{psmallmatrix}$，闭环从扰动 $\mathbf w$ 到 $\mathbf z$ 的传递 $\mathbf T_{zw}(s)=\begin{psmallmatrix}\mathbf Q^{1/2}\\-\mathbf R^{1/2}\mathbf K\end{psmallmatrix}(s\mathbf I-\mathbf A_{cl}^\star)^{-1}$。则 LQR 最优代价恰是 $H_2$ 范数平方：
\begin{equation}
  J^\star=\tfrac12\mathbf x_0^\top\mathbf P^\star\mathbf x_0=\tfrac12\|\mathbf T_{zw}\|_2^2\cdot\|\mathbf x_0\|^2.
  \label{eq:lqr-h2}
\end{equation}
即 LQR $=$ $H_2$ 最优控制——能量意义下最好的线性反馈。从 $H_2$ 到 $H_\infty$ 只需把“能量最优”改为“最坏情况最优”。

\subsection{\texorpdfstring{$H_\infty$}{H-infinity} 即 minimax 博弈与博弈型 Riccati}
\label{sec:lqr-hinf-game}

设扰动经 $\dot{\mathbf x}=\mathbf A\mathbf x+\mathbf B\mathbf u+\mathbf D\mathbf w$ 进入。$H_\infty$ 把控制问题变成\emph{零和微分博弈}：控制器 $\mathbf u$ 最小化代价，自然界的“对手”$\mathbf w$ 最大化代价：
\begin{equation}
  \min_{\mathbf u}\max_{\mathbf w}\int_0^\infty\big(\|\mathbf z\|^2-\gamma^2\|\mathbf w\|^2\big)\mathrm dt.
  \label{eq:lqr-hinf-game}
\end{equation}
$\gamma$ 是\emph{扰动衰减水平}：设计者承诺，无论对手 $\mathbf w$ 怎样使坏，评价输出 $\mathbf z$ 的能量不超过 $\gamma^2$ 倍扰动能量。$\gamma$ 越小、承诺越强、越鲁棒，但越保守（控制能量越大）。

\begin{theorem}[博弈型 ARE\cite{zhou1996robust}]\label{thm:lqr-game-are}
设 $V(\mathbf x)=\tfrac12\mathbf x^\top\mathbf P\mathbf x$。对博弈 \cref{eq:lqr-hinf-game}，Isaacs（博弈 HJB）方程
\begin{equation}
  0=\min_{\mathbf u}\max_{\mathbf w}\Big\{\tfrac12\mathbf x^\top\mathbf Q\mathbf x+\tfrac12\mathbf u^\top\mathbf R\mathbf u-\tfrac12\gamma^2\mathbf w^\top\mathbf w+\mathbf x^\top\mathbf P(\mathbf A\mathbf x+\mathbf B\mathbf u+\mathbf D\mathbf w)\Big\}
  \label{eq:lqr-isaacs}
\end{equation}
的鞍点为 $\mathbf u^\star=-\mathbf R^{-1}\mathbf B^\top\mathbf P\mathbf x$、$\mathbf w^\star=\gamma^{-2}\mathbf D^\top\mathbf P\mathbf x$，且 $\mathbf P$ 满足博弈型代数 Riccati 方程
\begin{equation}
  \mathbf A^\top\mathbf P+\mathbf P\mathbf A+\mathbf P\big(\gamma^{-2}\mathbf D\mathbf D^\top-\mathbf B\mathbf R^{-1}\mathbf B^\top\big)\mathbf P+\mathbf Q=\mathbf0.
  \label{eq:lqr-game-are}
\end{equation}
\end{theorem}

\begin{proof}[推导骨架]
对 \cref{eq:lqr-isaacs} 大括号关于 $\mathbf u$ 求极小（$\partial/\partial\mathbf u:\mathbf R\mathbf u+\mathbf B^\top\mathbf P\mathbf x=\mathbf0$）得 $\mathbf u^\star$；关于 $\mathbf w$ 求极大（$\partial/\partial\mathbf w:-\gamma^2\mathbf w+\mathbf D^\top\mathbf P\mathbf x=\mathbf0$，Hessian $-\gamma^2\mathbf I\prec\mathbf0$ 确认极大）得 $\mathbf w^\star$。代回并对所有 $\mathbf x$ 提取系数即 \cref{eq:lqr-game-are}。
\end{proof}

\begin{insight}
\cref{eq:lqr-game-are} 与 CARE 的差别只在中间一项的符号结构：$-\mathbf P\mathbf B\mathbf R^{-1}\mathbf B^\top\mathbf P$ 是控制的“好处”（减小代价），$+\gamma^{-2}\mathbf P\mathbf D\mathbf D^\top\mathbf P$ 是对手的“破坏”（增大代价）。两者竞争决定 ARE 是否有解：
\begin{itemize}
  \item $\gamma\to\infty$ 时 $\gamma^{-2}\mathbf D\mathbf D^\top\to\mathbf0$，博弈 ARE \emph{退化为} CARE \cref{eq:lqr-care}——$H_\infty$ 是 LQR 的“最坏情况推广”；
  \item $\gamma$ 减小，$+\gamma^{-2}(\cdot)$ 项使 Riccati 右端越来越大，直到临界 $\gamma^\star$ 处正定稳定化解\emph{消失}，对应“系统抗不住这么强的扰动”。有限时间版的 Riccati 在此发生有限时间爆破，对应博弈无解。
\end{itemize}
最小的可行 $\gamma^\star$ 即闭环 $\|\mathbf T_{zw}\|_\infty$；计算上用二分（$\gamma$-迭代），每步解一个 ARE 判可行。这与 RL 中 robust MDP 的 $\min_\pi\max_{P\in\mathcal P}$、域随机化是\emph{相同的数学结构}。
\end{insight}

\subsection{LTR：恢复 LQG 失去的裕度}
\label{sec:lqr-ltr}

回到 \cref{sec:lqr-doyle} 的难题：LQG 没有保证裕度，怎么补救？最常用的工程手段是 \emph{回路传递恢复}（Loop Transfer Recovery, LTR）：人为增大过程噪声协方差 $\boldsymbol\Sigma_w\leftarrow\boldsymbol\Sigma_{w,0}+q\,\mathbf B\mathbf B^\top$，令 $q\to\infty$。此时 Kalman 增益 $\mathbf L$ 在 $\mathbf B$ 的列空间方向变大，观测器极点被推到远离虚轴处，LQG 的回路传递渐近趋于 LQR 的回路传递——从而“恢复”LQR 的好裕度。其限制有三：仅对\emph{最小相位}系统有效（开环有右半平面零点时无法完全恢复）；增大虚拟 $\boldsymbol\Sigma_w$ 放大观测噪声、控制信号抖动；只恢复\emph{输出端}（回路断开点）裕度，非输入端。

\begin{remark}[分工声明]
本节只做 $H_\infty$ 的 Riccati/博弈入口（状态空间侧、与 LQR 的代数连续性）。\textbf{频域 $H_\infty$ 范数定义、灵敏度 $S/T$ 与回路成形、$\mu$-综合/DK 迭代、DGKF 两-Riccati 状态空间解（Doyle--Glover--Khargonekar--Francis 1989\cite{doyle1989dgkf}）全部见兄弟章 \cref{ch:robust-hinf}}——代数侧在本章、频域侧在 E 章。博弈的\emph{规划/多智能体}侧（iLQGames、HJI 可达）属规划部与兄弟章 \cref{ch:hj-reachability}，本章不重复 HJI/可达。
\end{remark}

% ════════════════════════════════════════════════════════════════
\section{ARE 数值求解：算法谱系与收敛性}
\label{sec:lqr-numerics}

理论上 CARE 解存在唯一（\cref{sec:lqr-existence}），但工程要的是\emph{数值}解。本节梳理 ARE 求解的算法谱系及其收敛性，并把数值代码剔除、只留理论与算法骨架。

\subsection{Schur 法：稳定不变子空间的数值构造}
\label{sec:lqr-schur}

Laub 1979\cite{laub1979schur} 提出用 Hamilton 矩阵的实 Schur 分解构造稳定不变子空间，是 MATLAB \texttt{care}/\allowbreak\texttt{icare}、SciPy \texttt{solve\_continuous\_are}、SLICOT 的默认算法。

\begin{algo}[Schur 法求解 CARE\cite{laub1979schur}]\label{algo:lqr-schur}
\begin{enumerate}
  \item 构造 Hamilton 矩阵 $\mathcal Z=\begin{psmallmatrix}\mathbf A&-\mathbf S\\-\mathbf Q&-\mathbf A^\top\end{psmallmatrix}$（$\mathbf S=\mathbf B\mathbf R^{-1}\mathbf B^\top$）；
  \item 实 Schur 分解 $\mathbf U^\top\mathcal Z\mathbf U=\mathbf T$（$\mathbf U$ 正交、$\mathbf T$ 准上三角），并\emph{排序}使左上 $n\times n$ 块对应全部 $\operatorname{Re}\lambda<0$ 的特征值（QR 迭代 $+$ Givens 旋转交换，LAPACK \texttt{dgees} 带排序选项一步完成）；
  \item 取 $\mathbf U$ 前 $n$ 列分块 $\begin{psmallmatrix}\mathbf U_{11}\\\mathbf U_{21}\end{psmallmatrix}$，它张成稳定不变子空间；
  \item 解 $\mathbf P^\star=\mathbf U_{21}\mathbf U_{11}^{-1}$，强制对称 $\mathbf P\leftarrow(\mathbf P+\mathbf P^\top)/2$；
  \item 验证残差 $\|\mathbf A^\top\mathbf P+\mathbf P\mathbf A-\mathbf P\mathbf S\mathbf P+\mathbf Q\|<\varepsilon$、闭环谱 $\subset\mathbb C_-$。
\end{enumerate}
\end{algo}

\begin{table}[H]\centering\small
\caption{Schur 法 vs 直接特征分解}
\label{tab:lqr-schur-vs-eig}
\begin{tabular}{lll}
\toprule
方面 & 特征分解 & Schur 分解 \\
\midrule
数值稳定性 & 特征向量矩阵可能病态（条件数大） & $\mathbf U$ 正交，条件数 $=1$ \\
复数运算 & 有复特征值时需复运算 & 实 Schur 全程实数 \\
对称性保持 & 不自动保证 $\mathbf P=\mathbf P^\top$ & 易后处理强制对称 \\
复杂度 & $O(n^3)$ & $O(n^3)$，常数更小 \\
\bottomrule
\end{tabular}
\end{table}

\subsection{Newton--Kleinman 迭代 \texorpdfstring{$=$}{=} 策略迭代}
\label{sec:lqr-newton-kleinman}

Schur 是“一步到位”的直接法，但对大规模系统（$n>100$）代价高。Newton--Kleinman 是\emph{策略迭代}：从稳定初始增益 $\mathbf K_0$ 出发，反复解 Lyapunov 方程（线性！），二次收敛到 CARE 解\cite{kleinman1968iterative}。

\begin{algo}[Newton--Kleinman 迭代\cite{kleinman1968iterative}]\label{algo:lqr-kleinman}
\begin{enumerate}
  \item \textbf{初始化}：选 $\mathbf K_0$ 使 $\mathbf A_0=\mathbf A-\mathbf B\mathbf K_0$ Hurwitz（极点配置或 LQR 粗解）；
  \item \textbf{策略评估}：解 Lyapunov 方程 $\mathbf A_k^\top\mathbf P_k+\mathbf P_k\mathbf A_k+\mathbf Q+\mathbf K_k^\top\mathbf R\mathbf K_k=\mathbf0$ 得 $\mathbf P_k$（$\mathbf A_k$ Hurwitz $+$ 右端正定 $\Rightarrow$ 唯一正定解）；
  \item \textbf{策略改进}：$\mathbf K_{k+1}=\mathbf R^{-1}\mathbf B^\top\mathbf P_k$；
  \item 重复至 $\|\mathbf P_{k+1}-\mathbf P_k\|<\varepsilon$。
\end{enumerate}
\end{algo}

\begin{proposition}[Kleinman 迭代 $=$ Riccati 算子的 Newton 法]\label{prop:lqr-kleinman-newton}
CARE 残差算子 $\mathcal R(\mathbf P)=\mathbf A^\top\mathbf P+\mathbf P\mathbf A-\mathbf P\mathbf B\mathbf R^{-1}\mathbf B^\top\mathbf P+\mathbf Q$ 在 $\mathbf P_k$ 处的 Fréchet 导数为 $\mathcal R'(\mathbf P_k)\boldsymbol\Delta=(\mathbf A-\mathbf B\mathbf K_k)^\top\boldsymbol\Delta+\boldsymbol\Delta(\mathbf A-\mathbf B\mathbf K_k)$。Newton 步 $\mathcal R'(\mathbf P_k)\boldsymbol\Delta=-\mathcal R(\mathbf P_k)$ 恰是 \cref{algo:lqr-kleinman} 第 2 步的 Lyapunov 方程（$\boldsymbol\Delta=\mathbf P_{k+1}-\mathbf P_k$）。
\end{proposition}

\begin{theorem}[收敛性，Kleinman 1968 / Hewer 1971\cite{kleinman1968iterative,hewer1971iterative}]\label{thm:lqr-kleinman-conv}
从任意使 $\mathbf A_0$ Hurwitz 的 $\mathbf K_0$ 出发，\cref{algo:lqr-kleinman} 满足：
\begin{enumerate}
  \item \textbf{单调性}：$\mathbf P_1\succeq\mathbf P_2\succeq\cdots\succeq\mathbf P^\star$，自\emph{第一个} Newton 迭代起从上方单调逼近\footnote{单调链由 $\mathbf P_1$ 而非初值 $\mathbf P_0$ 起算：任取一个使 $\mathbf A_0$ Hurwitz 的初始增益 $\mathbf K_0$，其闭环代价 $\mathbf P_0$ 并不必然支配 $\mathbf P^\star$（次优反馈在某些方向可能比最优更激进），只有自第一步迭代 $\mathbf P_1$ 起才有 $\mathbf P_1\succeq\mathbf P_2\succeq\cdots\succeq\mathbf P^\star$。此处指标依 Kleinman(1968)\cite{kleinman1968iterative} 原文，亦见 Riccati 迭代法综述对 CARE 牛顿法“自 $\mathbf X_1$ 从上方逼近”的标准陈述。}；
  \item \textbf{稳定性保持}：每个 $\mathbf A_k=\mathbf A-\mathbf B\mathbf K_k$ Hurwitz（不中途跳出稳定域）；
  \item \textbf{二次收敛}：$\|\mathbf P_{k+1}-\mathbf P^\star\|\le C\|\mathbf P_k-\mathbf P^\star\|^2$。
\end{enumerate}
\end{theorem}

单调性证明思路（对 $k\ge1$）：$V_k(\mathbf x)=\mathbf x^\top\mathbf P_k\mathbf x$ 是用反馈 $\mathbf u=-\mathbf K_k\mathbf x$ 的闭环代价（$\mathbf P_k$ 满足以 $\mathbf K_k$ 为反馈的 Lyapunov 方程）；$\mathbf K_{k+1}=\mathbf R^{-1}\mathbf B^\top\mathbf P_k$ 是在 $\mathbf P_k$ 基础上对 Bellman 右端做一步贪心改进，故 $V_{k+1}\le V_k$，即 $\mathbf P_{k+1}\preceq\mathbf P_k$。注意该比较自第一步迭代起生效——初值 $\mathbf K_0$ 仅需保证 $\mathbf A_0$ Hurwitz，并不要求其闭环代价 $\mathbf P_0$ 支配 $\mathbf P^\star$，故单调链由 $\mathbf P_1$ 起算。离散稳态增益的局部二次收敛由 Hewer 1971 给出。

\subsection{SDA、符号函数与广义 ARE}
\label{sec:lqr-sda-sign}

\paragraph{辛 Doubling（SDA）。}对大规模稀疏系统（如 MPC 内部多步 Riccati 递推），SDA 利用辛矩阵的“翻倍”性质：每次迭代等效把时间步长加倍，收敛步数 $O(\log(1/\varepsilon))$。它是 HPIPM（acados 内部 QP 求解器）使用的核心算法（库细节不展开）。

\paragraph{矩阵符号函数（Roberts 1980）。}对 Hamilton 矩阵 $\mathcal Z$ 定义符号函数 $\operatorname{sign}(\mathcal Z)=\mathcal Z(\mathcal Z^2)^{-1/2}$，其特征值为原谱的符号（$\pm1$）。迭代 $\mathcal Z_{k+1}=\tfrac12(\mathcal Z_k+\mathcal Z_k^{-1})$ 收敛到 $\operatorname{sign}(\mathcal Z)$，稳定投影 $\boldsymbol\Pi_-=\tfrac12(\mathbf I-\operatorname{sign}(\mathcal Z))$ 给出稳定子空间，进而算 $\mathbf P^\star$\cite{roberts1980linear}。优势：只用矩阵逆，适合并行/GPU。

\paragraph{广义 ARE。}对带代数约束的描述子系统 $\mathbf E\dot{\mathbf x}=\mathbf A\mathbf x+\mathbf B\mathbf u$（$\mathbf E$ 奇异），标准 CARE 不适用，需广义 ARE $\mathbf A^\top\mathbf X\mathbf E+\mathbf E^\top\mathbf X\mathbf A-\mathbf E^\top\mathbf X\mathbf B\mathbf R^{-1}\mathbf B^\top\mathbf X\mathbf E+\mathbf Q=\mathbf0$，用广义 Schur（QZ 算法）求解。应用：多体系统的 DAE 形式（如物理引擎内部）、电力系统微分--代数模型。

\subsection{值/策略迭代 \texorpdfstring{$=$}{=} RL 的精确对应}
\label{sec:lqr-rl-bridge}

ARE 数值方法在 RL 里有精确对应，这把经典控制与现代 RL 焊在一起，也是通往 P8 强化学习部的代数桥。

\begin{table}[H]\centering\small
\caption{Riccati 数值方法与 RL 概念的精确映射}
\label{tab:lqr-rl-map}
\begin{tabular}{lll}
\toprule
经典控制 & RL 对应 & 收敛阶 \\
\midrule
DARE 递推（值函数倒推） & 值迭代（Value Iteration） & 线性 \\
Newton--Kleinman（解 Lyapunov $+$ 改进增益） & 策略迭代（Policy Iteration） & 二次 \\
$\mathbf K=(\mathbf R+\mathbf B^\top\mathbf P\mathbf B)^{-1}\mathbf B^\top\mathbf P\mathbf A$（贪心，离散） & $\pi^\star=\arg\max_a Q(s,a)$ & --- \\
Fazel 梯度下降\cite{fazel2018global} & 策略梯度（REINFORCE） & 线性（PL 条件） \\
自然策略梯度 & Natural Policy Gradient & Fisher $=$ Gramian 归一化 \\
\bottomrule
\end{tabular}
\end{table}

\begin{insight}
值迭代 vs 策略迭代的差别，在 RL 中是 Q-learning vs Actor--Critic 的差别。LQR 是唯一能\emph{精确}分析两者收敛率差异的问题：策略迭代的二次收敛（$=$ Newton 法）解释了为何 Actor--Critic 在实践中常比纯 Q-learning 更快收敛。更妙的是，Fazel 等\cite{fazel2018global} 分析的三种更新（普通/自然策略梯度、Gauss--Newton）中，\emph{Gauss--Newton} 步（取步长 $\tfrac12$）恰等价于 Kleinman 的一步策略迭代——经典控制（1968）与现代 RL 在算法层面完美统一\footnote{此处须区分：与 Kleinman(1968)/Hewer(1971) 的策略迭代一步重合的是 Fazel 等的 \emph{Gauss--Newton}（步长 $\tfrac12$）更新，而非自然策略梯度（后者按 Fisher 信息度规预条件、是介于普通梯度与 Gauss--Newton 之间的另一种更新）。}。本表只作“代数桥”，RL 算法主线见 P8 强化学习运动控制部，不在此展开。
\end{insight}

% ════════════════════════════════════════════════════════════════
\section{前沿与出口}
\label{sec:lqr-frontier}

本章把 LQR/Riccati 的“代数主干”讲透了。它向四个方向延伸，这里点名而不展开，给出交叉引用与文献出口。

\paragraph{流形 LQR / TVLQR / iLQR。}状态空间是流形（如 $\SO(3)$、$\SE(3)$）时不能直接线性化，标准做法是\emph{切空间 LQR}：用指数坐标 $\delta\boldsymbol\phi\in\so(3)$、右扰动 $\mathbf J_r$（\cref{ch:lie}）把动力学在李代数里线性化 $\dot{\delta\boldsymbol\phi}=\mathbf A(t)\delta\boldsymbol\phi+\mathbf B(t)\delta\mathbf u$，再用常规时变 LQR。给定标称轨迹的\emph{时变} LQR（TVLQR）做轨迹镇定，其倒向 RDE 递推正是 iLQR backward pass 的精确形式——“iLQR 的 backward pass $=$ TVLQR $=$ 本章 Riccati 递推沿轨迹版”，TVLQR/iLQR/DDP 的展开见兄弟章 \cref{ch:ddp-ilqr}。

\paragraph{LQR 作为学习型控制的果蝇。}Fazel 等 2018\cite{fazel2018global} 证明 LQR 代价 $C(\mathbf K)$ 虽\emph{非凸非拟凸}，却满足梯度优势（PL 条件），（自然/投影）策略梯度\emph{线性全局收敛}、无虚假局部极小——这是第一个非凸连续控制问题被严格证明全局收敛的案例，接 P8 RL 部。其余前沿各一句：样本复杂度（Coarse-ID $+$ SLS 给 $\tilde O((n+m)/\varepsilon^2)$ 样本，Dean 等 2020\cite{dean2020sample}）；在线 LQ 悔界（首个 $\tilde O(\sqrt T)$，Abbasi-Yadkori--Szepesvári 2011\cite{abbasiyadkori2011regret}）；分布鲁棒 LQ（噪声落在以高斯为心的 Wasserstein 球内，最优策略仍为线性观测反馈，Taşkesen 等 2023\cite{taskesen2023distributionally}）；平均场 LQG（$N\to\infty$ 弱耦合代理的去中心化 $\varepsilon$-Nash，Huang--Malhamé--Caines 2006\cite{huang2006large}）。

\paragraph{Transformer 与 in-context Kalman。}Goel--Bartlett 2024\cite{goel2024transformer} 显式构造权重，证明因果 softmax self-attention 可一致近似任意可观 LTI 系统的 Kalman 滤波器，并推广到 LQG；其理论基石是 Transformer 在 in-context learning 中实现梯度下降/闭式岭回归\cite{garg2022what,akyurek2023what}。

\begin{pitfall}[勿把“in-context 线性回归”错记为 LQR 结果]
Garg 等 2022\cite{garg2022what} 的“ICL 作为线性回归”是关于线性回归/决策树/两层网的案例研究，\emph{并不直接针对 LQR}。真正把 Transformer 接到 Kalman/LQG 的是 Goel--Bartlett 2024\cite{goel2024transformer}——引用时应以后者为主、前者作铺垫，切勿把“线性回归案例”错记到 LQR 头上。
\end{pitfall}

\paragraph{LMI 形式。}CARE 可经 Schur 补等价写成线性矩阵不等式 $\begin{psmallmatrix}\mathbf A^\top\mathbf P+\mathbf P\mathbf A+\mathbf Q&\mathbf P\mathbf B\\\mathbf B^\top\mathbf P&\mathbf R\end{psmallmatrix}\succeq\mathbf0$、$\mathbf P\succeq\mathbf0$，用 SDP 求解器统一处理（\cref{ch:nlopt} 的 SDP/凸优化）。LMI 不是 CARE 的“简化版”，而是\emph{更一般}的表示——它能加入参数不确定性（鲁棒 LQR）与结构约束（分散控制）这些 CARE 处理不了的问题，CARE 是 LMI 在确定性无约束情形下的特化。

\begin{remark}[cheap-control 渐近：LQR 极点的 Butterworth 模式]
有一条常被无条件化叙述的结论：“高阶系统的 LQR 极点趋于 Butterworth 分布”。精确陈述是 \emph{cheap-control 渐近}——当控制权重 $r\to0$（对称根轨迹 $1+\rho\,P(-s)P(s)=0$，$\rho\to\infty$）时，部分闭环极点沿 Butterworth 模式趋于无穷、其余趋于被控对象的（镜像）零点；这\emph{并非}对任意 $\mathbf Q,\mathbf R$ 普遍成立。仅在此 $r\to0$/对称根轨迹语境下才宜采用此说法。
\end{remark}

% ════════════════════════════════════════════════════════════════
\section*{常见误解汇总}
\addcontentsline{toc}{section}{常见误解汇总}

\begin{enumerate}
  \item \textbf{“CARE 有解就唯一”}：错。CARE 一般多解；\emph{唯一}的是\emph{稳定化半正定}解（\cref{thm:lqr-care-exist}），且它是最大半正定解（\cref{thm:lqr-maximal}）。数值上必须锁定\emph{稳定}不变子空间。
  \item \textbf{“LQG 是默认稳健方案”}：错。分离原理保证性能最优，不保证鲁棒裕度；Doyle 1978（\cref{ex:lqr-doyle}）证 LQG 无保证裕度。有参数不确定时须上 $H_\infty$ 或 LTR。
  \item \textbf{“LQR 好裕度对参数变化也成立”}：错。裕度是对\emph{回路增益}扰动、且要求\emph{状态可测、无饱和}（\cref{sec:lqr-disk-margin} 三前提）；对模型参数变化无直接保证。
  \item \textbf{“分离原理普遍成立”}：错。仅线性--高斯--二次下严格（\cref{tab:lqr-sep-fail}）；EKF $+$ LQR 是\emph{近似}分离，剧烈运动下可能崩。
  \item \textbf{“多输入 LQR 裕度和单输入一样好”}：不一定。矩阵回差不等式仍成立，但\emph{单通道} disk margin 可 $<\tfrac12$（\cref{sec:lqr-disk-margin}）。
  \item \textbf{“博弈 ARE 总有正定解”}：错。$\gamma<\gamma^\star$ 时正定稳定化解消失（\cref{thm:lqr-game-are} 后注），对应系统抗不住该扰动、博弈无解。
  \item \textbf{“DARE 中 $\mathbf R+\mathbf B^\top\mathbf P\mathbf B$ 的 $\mathbf P$ 取本步下标”}：错。应取 $k+1$（后一时刻），手推常错（呼应 \cref{ch:control} 离散陷阱）。
\end{enumerate}

% ════════════════════════════════════════════════════════════════
\section*{小结}
\addcontentsline{toc}{section}{小结}

本章在 \cref{ch:control} 与 C0 子部之上，把 LQR/LQG/Riccati 的四件“陈述但未证”之事补成完整证明，再向上接 $H_\infty$、数值谱系与学习。下表只列\emph{本章新增}的增量（导论已有的 RDE/CARE/DARE/裕度结论等不重列）。

\begin{table}[H]\centering\small
\caption{本章新增结论速查（深化于 \cref{ch:control}）}
\label{tab:lqr-deep-summary}
\begin{tabular}{p{0.30\linewidth}p{0.40\linewidth}p{0.22\linewidth}}
\toprule
增量 & 核心结果 & 出处 \\
\midrule
CARE 存在唯一性 & 可稳定化 $+$ 可检测 $\Rightarrow$ 唯一稳定化 $\mathbf P^\star\succeq\mathbf0$、闭环 Hurwitz；$\mathbf P^\star=$ 最大半正定解 & \cref{thm:lqr-care-exist,thm:lqr-maximal} \\
辛几何 why & $\mathbf P^\star=\mathbf X_2\mathbf X_1^{-1}$；$\boldsymbol\lambda=\mathbf P\mathbf x$ $=$ 稳定 Lagrangian 子空间的图 & \cref{thm:lqr-invsubspace} \\
回差不等式 & $\mathbf F^*\mathbf R\mathbf F\succeq\mathbf R$ $\Rightarrow$ 单输入 $\mathrm{GM}\in[\tfrac12,\infty)$、$\mathrm{PM}\ge60^\circ$ & \cref{thm:lqr-kalman-ineq,thm:lqr-simargin} \\
LQG 失裕机理 & Kalman $=$ 开环观测器，噪声沿不可观方向 $\Rightarrow$ 裕度坍缩 & \cref{ex:lqr-doyle} \\
分离原理 & 正交分解 $+$ 误差 Riccati 不含 $\mathbf u$ $\Rightarrow$ 估计/控制独立设计 & \cref{thm:lqr-separation} \\
$H_\infty$ 博弈 Riccati & $\mathbf A^\top\mathbf P+\mathbf P\mathbf A+\mathbf P(\gamma^{-2}\mathbf D\mathbf D^\top-\mathbf B\mathbf R^{-1}\mathbf B^\top)\mathbf P+\mathbf Q=\mathbf0$；$\gamma\to\infty$ 退化 CARE & \cref{thm:lqr-game-are} \\
数值谱系 & Schur（直接）/ Newton--Kleinman（策略迭代、二次收敛）/ SDA / 符号函数 & \cref{algo:lqr-schur,thm:lqr-kleinman-conv} \\
RL 桥 & 值迭代 $=$ DARE 递推、策略迭代 $=$ Newton--Kleinman & \cref{tab:lqr-rl-map} \\
\bottomrule
\end{tabular}
\end{table}

\paragraph{启下。}沿本章打开的三个接口走下去：博弈型 Riccati 的\emph{频域}另一半（$H_\infty$ 范数、灵敏度、$\mu$/DGKF）在 \cref{ch:robust-hinf}；把 Riccati 递推沿\emph{非线性轨迹}展开即 iLQR/DDP（\cref{ch:ddp-ilqr}）；本章用到的 Lyapunov/LaSalle 工具的系统化在 \cref{ch:lyapunov-stability}；而值/策略迭代这座“代数桥”通向 P8 强化学习运动控制部——LQR 正是经典控制跨入学习型机器人控制的主干道路。

% ════════════════════════════════════════════════════════════════
\section*{延伸阅读}
\addcontentsline{toc}{section}{延伸阅读}

\begin{itemize}
  \item \textbf{存在唯一性与 ARE 代数理论}：Anderson \& Moore, \emph{Optimal Control: Linear Quadratic Methods}\cite{anderson2007optimal}（LQR/LQG 工程圣经，含 LTR）；Lancaster \& Rodman, \emph{Algebraic Riccati Equations}\cite{lancaster1995algebraic}（ARE 代数理论最完整专著）；Kalman 1960\cite{kalman1960}、Wonham 1968\cite{wonham1968riccati} 原典。
  \item \textbf{鲁棒裕度与 $H_\infty$}：Kalman 1964\cite{kalman1964optimal}（回差不等式）；Safonov \& Athans 1977\cite{safonov1977gain}（多输入）；Doyle 1978\cite{doyle1978guaranteed}（“There are none.”）；Zhou, Doyle \& Glover, \emph{Robust and Optimal Control}\cite{zhou1996robust} 与 DGKF 1989\cite{doyle1989dgkf}（$H_\infty$ 圣经与状态空间解，频域细节见 \cref{ch:robust-hinf}）。
  \item \textbf{数值方法}：Laub 1979\cite{laub1979schur}（Schur）；Kleinman 1968\cite{kleinman1968iterative} 与 Hewer 1971\cite{hewer1971iterative}（Newton 迭代与收敛）；Roberts 1980\cite{roberts1980linear}（符号函数）。
  \item \textbf{分离原理与 LQG}：Kwakernaak \& Sivan, \emph{Linear Optimal Control Systems}\cite{kwakernaak1972linear}。
  \item \textbf{学习型 LQR（前沿）}：Fazel 等 2018\cite{fazel2018global}（PG 全局收敛）；Dean 等 2020\cite{dean2020sample}（样本复杂度）；Abbasi-Yadkori \& Szepesvári 2011\cite{abbasiyadkori2011regret}（在线悔界）；Taşkesen 等 2023\cite{taskesen2023distributionally}（分布鲁棒）；Goel \& Bartlett 2024\cite{goel2024transformer}（Transformer 表示 Kalman）。
\end{itemize}
